\section{Liquidity, Price Pressure, and Fire Sales}
\label{sec:liquidity_fire_sales}

The previous section established that institutional positions and trading can be correlated across investors, and that such correlation is not, by itself, a source of risk. Correlated positioning becomes economically relevant only through a mechanism that translates common exposure into price movements unrelated to fundamentals. The literature on liquidity-driven trading and fire sales provides this mechanism.

\textcite{coval2007asset} show that mutual funds experiencing large outflows tend to reduce their existing positions for reasons unrelated to changes in the underlying fundamentals of the securities involved. When several distressed funds hold the same securities, their flow-driven sales can generate concentrated, non-fundamental selling pressure in those common holdings. Prices consequently decline not because information about future cash flows has changed, but because of a temporary imbalance between the demand for liquidity and the market's capacity to absorb the sales. They further document that these price effects are partially reversed over subsequent quarters, consistent with an interpretation in which the initial decline reflects price pressure rather than a permanent revision of fundamental value.

This mechanism has two implications that are central to the present thesis. First, it identifies overlapping institutional ownership, rather than aggregate institutional ownership alone, as a channel through which liquidity shocks are transmitted to securities: a stock is exposed to fire-sale risk to the extent that several of its owners may face correlated redemptions that force them to sell. Second, it implies that the relevant risk is asymmetric. Because the fire-sale mechanism considered here originates from forced selling in response to redemptions, its price impact predominantly affects the downside of the return distribution rather than being symmetric around zero. This asymmetry motivates the thesis's focus on downside risk, rather than on volatility more broadly, as the prediction target.